\section{工具使用型科研智能体的设计与审计}
\label{sec:agent-system}

\subsection{模型决策、数值行动与证据的分离}

系统将研究规划与定位求解分为两个不同层次。语言模型接收目标、方法登记表、近期实验观察、持久记忆和审阅意见，自行提出下一步研究行动；外层控制器解析并校验响应，调用登记的确定性工具，再把工具结果送回下一次模型请求。数值结果来自实际工具执行，不能由模型响应中的文字代替。算法\ref{alg:research-loop}给出这一闭环的结构。

用户提供的Creative-Intelligence项目\cite{creativeIntelligence}为研究流程与状态管理的分工提供了架构参考。本项目重新实现控制器、任务专用数学工具与评测协议，没有复制该仓库源码，也不把自身称为Moonshine系统的分支。

与固定候选列表的自动调参相比，本系统允许模型同时提出可证伪假设、解释预计发生的现象、登记一个已实现算子的组合，并根据失败结果修改方法。其自主性仍然有明确边界：模型不能在这个协议中任意生成并执行新的Python程序，也不能增添未经实现的求解器。所谓方法构造，是在静态、恒速与交叉验证门控、距离或对数残差、尺度形式、稳健损失、截尾和子集初始化等已实现机制之间组合。数学创新程度因而需要由具体研究成果判断，不能从“使用了语言模型”直接推出。

\begin{figure}[H]
\centering
\begin{tikzpicture}[box/.style={draw=navy!50,fill=paleblue,rounded corners=2pt,align=center,text width=37mm,minimum height=13mm,font=\small},arr/.style={-{Stealth[length=2mm]},draw=teal,thick},node distance=9mm and 12mm]
\node[box] (plan) {大模型规划器\\形成假设、选择动作};
\node[box,right=of plan] (control) {持久化控制器\\校验、预算、派发};
\node[box,right=of control] (tool) {研究工具与求解器\\实验、诊断、数学见证};
\node[box,below=of tool] (evidence) {真实实验结果\\逐例数据与诊断};
\node[box,below=of control] (memory) {记忆与审阅反馈\\保留失败，修订假设};
\node[box,below=of plan] (freeze) {方案冻结\\独立测试与校准};
\draw[arr] (plan)--(control);
\draw[arr] (control)--(tool);
\draw[arr] (tool)--(evidence);
\draw[arr] (evidence)--(memory);
\draw[arr] (memory.north)--++(0,4mm)-|([xshift=-4mm]plan.south);
\draw[arr] ([xshift=5mm]plan.south)--([xshift=5mm]freeze.north);
\end{tikzpicture}
\caption{智能体的研究闭环。动作来自实际模型响应，证据来自数值工具；冻结后才开放最终评测。}
\label{fig:architecture}
\end{figure}


\begin{table}[H]
  \centering
  \small
  \caption{运行时各组件的职责与证据边界。}
  \label{tab:agent-components}
  \begin{tabular}{p{0.18\linewidth}p{0.35\linewidth}p{0.36\linewidth}}
    \toprule
    组件 & 主要职责 & 输出和限制\\
    \midrule
    模型提供者 & 发出真实CLI请求，或显式回放既有响应 & 原请求、JSON响应、事件与用量；回放不算新推理\\
    研究控制器 & 校验动作、检查预算、派发工具、维护状态 & 事件链与状态快照；模型失败不切换成规则提案\\
    研究工具箱 & 登记方法、执行实验、诊断、消融、冻结 & 指标与逐例证据；开发接口不能选择测试或校准集\\
    定位模块 & 根据观测求解声明的方法组合 & 位置、诊断、工作量；不读取真实位置和污染标签\\
    审阅者调用 & 审查泄漏、对照、计算预算和结论范围 & 同模型的独立提示调用；不等于独立复现\\
    数学核验工具 & 执行指定反例及有限数值见证 & 数值支持或失败；不替代命题证明\\
    \bottomrule
  \end{tabular}
\end{table}

\begin{algorithm}[H]
  \small
  \DontPrintSemicolon
  \caption{带结构校验、证据记录和冻结边界的科研闭环}
  \label{alg:research-loop}
  \KwData{研究目标、登记工具、开发数据接口、计算预算与模型提供者}
  \KwResult{具有验证证据的冻结方法，或明确记录的停止状态}
  初始化或核验恢复状态、事件链和运行身份摘要\;
  \While{状态为运行中且尚未冻结}{
    检查人工停止标记及动作、实验、模型调用预算\;
    \If{达到预设的成功常规实验审阅间隔}{
      发起独立审阅者模型调用，保存其意见并加入记忆\;
    }
    将工具定义、登记方法、近期观察和持久记忆交给规划模型\;
    校验响应字段、动作枚举和参数JSON，保存原始模型响应\;
    \If{请求失败或结构非法且有限重试已耗尽}{
      记录提供者失败，不生成替代研究决策\;
      \Return{失败状态及已经保存的证据}\;
    }
    写入待执行工具标记和调用事件，再执行登记的研究工具\;
    保存工具结果或错误观察，更新证据索引与记忆\;
    \If{冻结动作通过工具箱的证据门槛}{
      保存方法及登记表摘要，关闭后续开发动作\;
    }
  }
  返回冻结状态或失败、预算耗尽、人工停止等实际状态\;
\end{algorithm}

\subsection{结构化决策与可组合工具}

规划器的响应包含公开决策摘要、研究假设、动作名称、参数JSON、预期观察和记忆备注。动作只能取自登记集合。控制器要求所有必填字段恰好存在、字段类型正确、动作合法且参数可解析为JSON对象，并拒绝非有限JSON数值。参数进入工具后，还由工具签名与定位配置验证再次约束。这里保存的是供研究者审阅的简洁决策理由，不是要求模型暴露内部思维链。

\texttt{read\_sources}返回研究问题与实现数学说明。来源文字属于证据材料，不拥有改变运行规则的权限。\texttt{register\_method}接收方法标识、假设、数学定义和声明式配置；同一运行中既有标识不能覆盖，因此修改后的方法具有新的版本标识。登记时用一个无关的小案例检查接口是否能运行，这只是接口探测，不计入研究表现。

\texttt{run\_experiment}在相同案例上评价指定方法，返回总体与分场景误差、尾部指标、工作量及证据路径。开发期只有发现集和验证集接口，验证调用还有独立次数限制。\texttt{diagnose}在已执行结果中读取失败案例、模型选择频率与求解诊断；其看到的真实误差来自离线评价端，而求解器只接触观测字段。这一区分允许智能体研究开发集上的失败，却不把未知案例真值作为部署输入。

\texttt{ablate}复制已登记配置并施加模型提出的组件修改，然后在共同的发现集上比较两者。该工具保存改动内容，但因果解释仍需检查是否只改变了待研究因素；若一次同时修改多个字段，就不能归因于单一机制。\texttt{verify\_math}执行预先实现的数值见证。它检验给定推导是否出现数值矛盾，不自动证明任意新命题。

\texttt{freeze}将选择的方法、科学结论、验证结果和登记表摘要写入冻结文件。工具箱检查候选是否具有验证证据，以及是否完成最低限度的实验、多方法、诊断和数学核验。通过后开发工具关闭，独立测试与校准位于研究控制器之外。冻结门槛保证存在可审查过程，不保证选中了最优方法；负面或无显著改善的结论同样可能是合法研究结果。

\subsection{审阅者、记忆与方法修订}

审阅者是与规划步骤分开的模型调用，使用专门提示审查数据泄漏、缺少对照、计算预算不公平、未经证明的创新声明和遗漏的反例。它使用相同的配置模型，而不是一个独立训练的第二模型；“独立”仅指调用上下文和职责分开。它可能与规划器共享盲点，其意见本身既不是证明，也不构成独立实验复现。

审阅意见与每步假设、预期观察、执行状态一起进入持久记忆；下一次规划上下文还包含方法登记表、最近的工具观察和剩余预算。模型可以据此提出新的组合、修改参数、安排消融或停止某条思路。历史结果不会被替换成一段只保留成功的最终叙述：完整调用和失败均留在事件记录中。当前控制器按成功的常规实验调用触发周期审阅；消融工具虽然消耗实验预算，但不能把它的存在误报为一次额外审阅者推理。

这种结构支持对研究过程的几个独立问题进行核查：提案是否先于对应结果；修订是否引用了当时已经存在的证据；某个失败是否导致后续行动改变；最终方法是否在测试访问之前冻结。它并不要求每轮都改善结果。若模型只是反复登记近似配置，日志同样会暴露这一点，不能因为系统具有“记忆”和“审阅者”组件就断言其研究有效。

\subsection{持久化、故障与审计强度}

真实请求保存提示文本、输出结构、响应JSON、CLI原始事件、标准错误和元数据。元数据记录提供者、请求摘要、耗时以及能够获得的用量。控制器另外维护追加式事件记录，每条事件包含序号、上一事件摘要和自身规范JSON摘要；状态快照采用临时文件写入后替换。此设计能够发现普通日志损坏、丢序或与保存摘要不一致的修改，但它不是数字签名或第三方时间戳，不能声称对能重写全部文件的操作者具有不可篡改保证。

恢复运行时，控制器核对研究目标、工具定义、配置、种子摘要、提供者和预算等运行身份信息，并验证事件链。工具执行之前先保存待执行标记；如果程序在工具完成与状态提交之间中断，重启会拒绝盲目重做这一副作用不确定的调用。该设计牺牲了一部分自动恢复便利，以避免重复登记、覆盖或把未完成结果当成已完成证据。

模型请求允许有限重试。超过限制、超时或提供者失败时，系统保存失败状态并停止相关流程，不悄悄退回UCB、随机提案、硬编码脚本或伪造模型响应。工具错误则成为下一轮可见的观察，由模型决定是否修正后再调用。总动作数、实验调用数、模型调用数和人工停止标记共同约束运行；预算耗尽本身不等同于完成研究，也不会授权尚未冻结的方案访问最终测试。

真实CLI调用请求只读沙箱，并在提示中禁止模型直接调用CLI内置的文件、命令或网络工具，由外层控制器执行研究工具。然而，提示不构成完全隔离的技术保证；提供者会保留原始CLI事件并标记意外工具行为。如果出现此类记录，就不能继续声称本次运行是完全无内置工具的纯模型接口。具体是否发生、模型版本和用量，应由本次运行证据说明。

\subsection{真实推理与显式回放}

回放模式顺序读取已经保存的模型响应，用于检查协议、工具调度和结果再计算。它在状态和请求审计中标记为回放，保留来源请求索引，并声明没有进行新的模型推理；既有响应耗尽即失败，不自动切换到在线调用或脚本答案。可选的严格模式要求提示摘要与原请求一致，非严格模式则记录不一致，不能把相近上下文下重用响应等同于模型在新证据下重新决策。

因此，回放成功只能支持“这组已保存决策能够被重新执行或检查”，不能单独证明语言模型今天仍会选择相同行动，也不能替代真实运行的模型事件。本文的实证部分将模型产生的新决策、确定性数值计算、数值反例验证与回放审计作为不同证据类别。这样的区分使研究深度能够由实际行为和结果判断，而不依赖未经验证的系统标签。
